\subsection*{الف}
ابتدا تعداد کل پیکسل‌ها $N=100$ است، پس:
\begin{flalign*}
p_u(x_i) &= \left[ \tfrac{40}{100}, \tfrac{30}{100}, \tfrac{20}{100}, \tfrac{10}{100} \right]
= [0.4,\,0.3,\,0.2,\,0.1] \\
p_v(y_i) &= \left[ \tfrac{0}{100}, \tfrac{20}{100}, \tfrac{50}{100}, \tfrac{30}{100} \right]
= [0,\,0.2,\,0.5,\,0.3] &&
\end{flalign*}

\subsection*{ب}
\begin{flalign*}
F_u(x_i) &= [0.4,\,0.7,\,0.9,\,1.0] \\
F_v(y_i) &= [0,\,0.2,\,0.7,\,1.0] &&
\end{flalign*}

\subsection*{ج}
در روش گسسته، برای هر $u$ مقدار $s=F_u(u)$ را محاسبه کرده و کوچک‌ترین $v$ را می‌یابیم که
$F_v(v)\ge s$:
\begin{flalign*}
u=0:\ s=0.4 &\Rightarrow v=2 \quad (F_v=0.7)\\
u=1:\ s=0.7 &\Rightarrow v=2 \quad (F_v=0.7)\\
u=2:\ s=0.9 &\Rightarrow v=3 \quad (F_v=1.0)\\
u=3:\ s=1.0 &\Rightarrow v=3 \quad (F_v=1.0) &&
\end{flalign*}
پس نگاشت:
\begin{flalign*}
0\to 2,\quad 1\to 2,\quad 2\to 3,\quad 3\to 3 &&
\end{flalign*}

\subsection*{د}
وارون‌گیری \lr{CDF} ورودی فقط همسان‌سازی هیستوگرام (به توزیع یکنواخت) را می‌دهد، نه مشخصه‌سازی به توزیع دلخواه. علاوه بر این، \lr{CDF} در حالت گسسته پله‌ای و غیر وارون‌پذیر است و چند سطح خاکستری می‌توانند به یک مقدار نگاشت شوند؛ بنابراین باید از تطبیق با \lr{CDF} هدف استفاده کرد.